\chapter{Authorization: Roles and Permissions}

Knowing \emph{who} is making a request is only half the security story.
Knowing \emph{what that person is allowed to do} is the other half. A logged
in user and a logged-in admin both pass the \texttt{protect} middleware from
Chapter 20 -- but only one of them should be able to delete someone else's
task. That distinction is authorization.

\section{Authentication vs Authorization}

These two terms get used interchangeably in casual conversation, but they
answer different questions and are implemented as separate middleware layers.

\begin{center}
\begin{tabularx}{\textwidth}{lXX}
\toprule
\textbf{} & \textbf{Authentication} & \textbf{Authorization} \\
\midrule
Question answered & "Who are you?" & "What are you allowed to do?" \\
When it runs & Every protected request, first & After authentication, before the controller \\
Implemented by & \texttt{protect} middleware (Ch. 18/20) & \texttt{authorize(...roles)} middleware (this chapter) \\
Failure response & 401 Unauthorized & 403 Forbidden \\
Based on & Valid JWT / session & \texttt{req.user.role} (or other permission data) \\
\bottomrule
\end{tabularx}
\end{center}

\begin{notebox}[NOTE]
A precise mental shortcut: \textbf{401} means "I don't know who you are" (or
your proof of identity is invalid). \textbf{403} means "I know exactly who
you are, and the answer is still no."
\end{notebox}

\section{The Role Field}

Recall the User schema from Chapter 13 -- the \texttt{role} field constrains
every user to one of three values:

\begin{lstlisting}[language=JavaScript]
// models/User.js (relevant field)
role: { type: String, enum: ["user", "manager", "admin"], default: "user" },
\end{lstlisting}

\begin{itemize}
  \item \textbf{user} -- can manage only their own tasks.
  \item \textbf{manager} -- can view and update tasks across their team.
  \item \textbf{admin} -- full access, including deleting any task or user.
\end{itemize}

\section{The \texttt{authorize} Middleware Factory}

Rather than writing a separate middleware for every role combination, write
one factory function that takes the allowed roles as arguments and returns
a middleware tailored to that route.

\begin{lstlisting}[language=JavaScript]
// middleware/authorize.js
export const authorize = (...allowedRoles) => {
  return (req, res, next) => {
    if (!req.user) {
      return res.status(401).json({ success: false, message: "Not authenticated" });
    }

    if (!allowedRoles.includes(req.user.role)) {
      return res.status(403).json({
        success: false,
        message: `Role '${req.user.role}' is not permitted to perform this action`,
      });
    }

    next();
  };
};
\end{lstlisting}

Because \texttt{authorize} always runs after \texttt{protect}, it can rely
on \texttt{req.user} already being set -- the \texttt{!req.user} check is
just a defensive guard in case the middleware order is ever misconfigured.

\section{Using It on Routes}

\begin{lstlisting}[language=JavaScript]
// routes/taskRoutes.js
import { protect } from "../middleware/authMiddleware.js";
import { authorize } from "../middleware/authorize.js";
import { deleteTask, getTasks } from "../controllers/taskController.js";

router.get("/tasks", protect, getTasks);

router.delete("/tasks/:id", protect, authorize("admin", "manager"), deleteTask);
\end{lstlisting}

Reading this route top to bottom: a request must first pass \texttt{protect}
(prove identity), then \texttt{authorize("admin", "manager")} (prove
permission), before \texttt{deleteTask} ever executes. A plain \texttt{user}
role gets a 403 and the delete never touches the database.

\begin{tipbox}[TIP]
Because \texttt{authorize} takes roles as arguments, the same middleware
function covers every permission combination in the app --
\texttt{authorize("admin")} for admin-only routes,
\texttt{authorize("admin", "manager")} for shared access, with no
duplicated logic.
\end{tipbox}

\begin{warnbox}[WARNING]
Role checks alone don't cover ownership. A \texttt{manager} passing
\texttt{authorize("manager")} still shouldn't necessarily be able to delete
a task belonging to a different team. For ownership-level rules, add an
explicit check inside the controller (e.g. comparing \texttt{task.user}
against \texttt{req.user.\_id}) in addition to the role middleware.
\end{warnbox}

\whatwhyhow
{An \texttt{authorize(...roles)} middleware factory that checks \texttt{req.user.role} against a list of permitted roles for a given route.}
{Authentication only proves identity; without a separate authorization check, any logged-in user -- regardless of role -- could call sensitive endpoints like deleting tasks or managing other users.}
{Chain it after \texttt{protect} in the route definition: \texttt{router.delete("/tasks/:id", protect, authorize("admin", "manager"), deleteTask)}, returning 403 for any role not in the allowed list.}
